% HD-F1-T35. GPT-6 Astra Pro, /root/transport; 22 November 2026 UTC.
\section[HD-F1-T35: effective genuine-spin rank bound]{HD-F1-T35: an effective bound for genuine spin layers and the exact remaining orthogonal quotient}
\label{sec:f1-genuine-spin-large-rank-exclusion}

For an original remaining T2-OTHER source with
orthogonal projective socle in odd defining
characteristic, every genuine spin layer has
\[
                  \boxed{\qquad r<10^{31}.\qquad}
 \tag{F1T35.1}
\]
Here the natural dimension is \(N=2r+1\) for type
\(B_r\), and \(N=2r\) for types \(D_r\) and
\({}^{2}D_r\). More precisely, for \(r\ge10^{31}\)
the actual universal central-cover map to the
original perfect layer kills the canonical
\(-1\in\operatorname{Spin}_N(q)\). It therefore
factors through the literal natural orthogonal
quotient \(\Omega_N(q)\). The descended map and
every original outer action are specified below.
The already retained field bound is \(q\le27\);
the proof of (F1T35.1) applies to every odd \(q\).

\paragraph{1. The actual cover and the canonical kernel element.}
Retain the original faithful absolutely irreducible
\(G\subset GL_K(V)\), \(m=\dim_KV\ge27\), scalar
centre \(Z\), and almost simple projective image
with socle \(S\). Let \(P\) be the inverse image
of \(S\) and \(L=P'\). The established maps of
T16 and T32 give
\[
          P=ZL,\qquad L=L',\qquad Z(L)=L\cap Z.
 \tag{F1T35.2}
\]
Let \(H=\operatorname{Spin}_{2r+1}(q)\) or
\(\operatorname{Spin}^{\epsilon}_{2r}(q)\), as
appropriate. The full exceptional-cover table in
Malle--Testerman, Remark~24.19/Table~24.3, has
orthogonal rank at most four. In the present range
\(\pi:H\to S\) is the universal central cover.
Thus there is a unique homomorphism
\[
 \beta:H\longrightarrow L,
 \qquad (L\longrightarrow S)\beta=\pi.
 \tag{F1T35.3}
\]
Its image \(I\) satisfies \(L=Z(L)I\), and
perfectness gives \(L=[I,I]\subset I\), so it is
surjective. The same universal property gives the
unique lift of each original conjugation:
\[
 \widetilde a_g\in\operatorname{Aut}(H),
 \qquad \beta\widetilde a_g=c_g\beta,
 \qquad \pi\widetilde a_g=a_g\pi.
 \tag{F1T35.4}
\]
Existence and uniqueness are obtained by applying
the universal property to the central extension
whose projection is composed with \(a_g^{-1}\);
the inverse lift proves that the map is an
automorphism, exactly as in T32.

The natural Clifford action is the exact sequence
\[
 1\longrightarrow\langle z_0\rangle
 \longrightarrow H\xrightarrow{\ \phi\ }
       \Omega_N(q)\longrightarrow1,
 \qquad z_0=-1.
 \tag{F1T35.5}
\]
This is the sequence used explicitly in the
published proof of Guralnick--Larsen--Tiep,
\emph{Inventiones mathematicae} \textbf{235}
(2024), 151--210, Proposition~6.3. We now verify
that all actual lifted automorphisms preserve its
displayed kernel, above the triality range.

Use the standard Euclidean coordinates
\(e_1,\ldots,e_r\) for the root system. For
\(B_r\), take simple roots
\(e_i-e_{i+1}\) and \(e_r\). For \(D_r\), take
\(e_i-e_{i+1}\) and the last root
\(e_{r-1}+e_r\). Write \(h_\alpha(t)=\alpha^\vee(t)\)
for the actual coroot map into a maximal torus
of the simply connected algebraic group. Then
\[
 z_0=
 \begin{cases}
 h_{\alpha_r}(-1),&B_r,\\
 h_{\alpha_{r-1}}(-1)h_{\alpha_r}(-1),&D_r.
 \end{cases}
 \tag{F1T35.6}
\]
Here is a direct verification of the element in
this formula. In both cases the cocharacter
lattice is the lattice of integral vectors with
even coordinate sum: the displayed simple
coroots generate it. Explicitly, for a vector
\(v=(v_1,\ldots,v_r)\) of even sum \(s\), its
coefficients in the \(B_r\) coroot basis are
\(v_1+\cdots+v_i\) for \(i<r\), and \(s/2\)
for \(i=r\). In the \(D_r\) basis they are
\(v_1+\cdots+v_i\) for \(i\le r-2\), followed
by \(s/2-v_r,s/2\). Substitution gives \(v\)
in each case. The two cocharacters in
(F1T35.6) are respectively \(2e_r\) and
\(2e_{r-1}\). Each belongs to this lattice and
does not belong to twice this lattice. Evaluation
at \(-1\) in a torus is trivial precisely when
all cocharacter coordinates are even in a lattice
basis; hence these are nonidentity elements of
order two. Every natural vector weight
\(\pm e_i\), and also zero in type \(B_r\),
pairs evenly with the indicated cocharacter.
Their natural images are therefore the identity.
The exact two-element kernel (F1T35.5) identifies
both displayed elements with the Clifford scalar
\(-1\).

The automorphism generators of Malle--Testerman
Theorem~24.24 and Steinberg Theorem~36 now fix
this element explicitly. Inner and diagonal
automorphisms lift to conjugations in the
connected algebraic group and fix its centre:
choose a preimage of the conjugating adjoint
point under the central isogeny over the
algebraic closure, and conjugate by that lift.
Field maps send \(h_\alpha(-1)\) to
\(h_\alpha((-1)^p)=h_\alpha(-1)\), since \(p\)
is odd. Type \(B_r\), \(r\ge3\), has no
nontrivial diagram symmetry. For \(D_r\),
\(r\ge5\), the only possible diagram operation
exchanges the last two simple roots, and hence
exchanges the two factors in (F1T35.6). The
twisted Frobenius in \({}^{2}D_r\) uses this same
exchange, so the displayed element is fixed and
lies in the finite group. The lifted field and
diagonal operations in the twisted case fix it
for the same formulas. The exceptional triality
diagram \(D_4\) lies below the asserted range.
By uniqueness in (F1T35.4), these standard lifts
are the actual lifts, and
\[
                    \widetilde a_g(z_0)=z_0
                         \quad\text{for every }g\in G.
 \tag{F1T35.7}
\]

\paragraph{2. The scalar on every actual constituent and the degree bound.}
Since \(\beta(z_0)\in Z(L)\subset Z\), it is an
actual scalar on \(V\), equal to either one or
minus one. Suppose it is minus one. In the actual
inflated decomposition
\[
              V_{\mathbf C}|_H=\bigoplus_i a_i W_i,
              \qquad m=\sum_i a_i\dim W_i,
 \tag{F1T35.8}
\]
the element \(z_0\) acts as minus one on every
\(W_i\). This retains all multiplicities; no
isotypy or extension of a chosen constituent is
assumed. In even natural dimension no \(W_i\)
factors through \(\phi\). In odd natural
dimension each \(W_i\) is faithful: \(H\) is
quasisimple with centre \(\langle z_0\rangle\)
(Malle--Testerman Corollary~24.13/Table~24.2),
and any proper normal kernel lies in that centre.
For clarity, a normal subgroup mapping onto the
simple quotient would give \(H=NZ(H)\), and
perfectness would force \(H=[N,N]\subset N\).
Thus a nontrivial character whose central
involution acts nontrivially has trivial kernel.

The exact degree bounds of the published
Proposition~6.3 cited above give
\[
 \begin{aligned}
 \dim W_i&>q^{r(r+1)/2}/4&& (N=2r+1),\\
 \dim W_i&>q^{r(r-1)/2}/4&& (N=2r).
 \end{aligned}
 \tag{F1T35.9}
\]
The full supplied proof was inspected: its odd
case uses the dual conformal-symplectic
semisimple index, and its even case restricts
along the actual Clifford embedding
\(\operatorname{Spin}_{2r-1}(q)\hookrightarrow
\operatorname{Spin}_{2r}^{\epsilon}(q)\), preserving
the same scalar \(-1\). These are named published
degree inputs, with their precise source scope
recorded separately.

\paragraph{3. All noncentral character ratios and the original commutator.}
Larsen--Tiep, \emph{Annals of Mathematics}
\textbf{200} (2024), 1--70, Theorem~5.5, applies
to these finite spin groups with natural
dimension \(N\). At its explicit constant
\[
 \sigma=\frac1{6507\cdot2^{21}\cdot10^{18}}
                >\sigma_0:=10^{-29},
 \qquad
 \frac{|\chi(h)|}{\chi(1)}
        \le\chi(1)^{-\sigma\operatorname{supp}(h)/N}.
 \tag{F1T35.10}
\]
Here support is measured on the actual natural
image \(\phi(h)\). Every noncentral \(h\in H\)
has support at least one. Indeed, support zero
makes \(\phi(h)\) scalar. For every \(x\in H\),
\([h,x]\) then belongs to the central kernel
\(\langle z_0\rangle\), so
\(x\mapsto[h,x]\) is a homomorphism to that
abelian group. Perfectness makes it trivial,
contradicting noncentrality of \(h\).

For \(r\ge10\), (F1T35.9), \(q\ge3\), and
\(N\le3r\) imply
\[
 \frac{\log_2\dim W_i}{N}
  >\frac{r-1}{6}-\frac{2}{3r}\ge\frac r7.
 \tag{F1T35.11}
\]
The last inequality, after multiplication by
\(42r\), is \(r^2-7r-28\ge0\), true at ten
and strictly increasing thereafter. At
\(r\ge10^{31}\),
\(\sigma_0r/7\ge100/7>4\).
Thus (F1T35.10) gives, for every noncentral
\(h\in H\),
\[
 \frac{|\chi_{W_i}(h)|}{\dim W_i}<\frac1{16},
 \qquad
 \left|\operatorname{Tr}_{V}\beta(h)\right|<\frac m{16}.
 \tag{F1T35.12}
\]
The second inequality is the sum over all actual
multiplicities in (F1T35.8).

Take an original nonscalar \(g\in G\). The
centralizer of \(S\) in its almost simple
automorphism group is trivial. Hence its induced
automorphism is nonidentity, and there is
\(e\in H\) for which
\[
 h_g=\widetilde a_g(e)e^{-1},\qquad
 \beta(h_g)=[g,\beta(e)]\notin Z(L).
 \tag{F1T35.13}
\]
In particular \(h_g\) is noncentral in \(H\).
After averaging a Hermitian form over the
original finite group, every original matrix
is unitary. The elementary inequality
\(\operatorname{Re}\operatorname{Tr}(u)
\ge m-2\operatorname{rank}(u-1)\), applied to
\(u=\beta(h_g)\), yields
\(\operatorname{rank}(\beta(h_g)-1)>15m/32\).
For \(e_L=\beta(e)\), the literal matrix identity
\[
 [g,e_L]-1
  =\bigl((g-1)e_L-e_L(g-1)\bigr)g^{-1}e_L^{-1}
 \tag{F1T35.14}
\]
then gives
\[
                    \operatorname{rank}(g-1)>15m/64.
 \tag{F1T35.15}
\]
Every nonidentity scalar has rank \(m\), so the
same weaker non-strict bound covers every true
inertia. No scalar phase or original inertia
has been replaced.

The original branch tuple has \(b\ge12\) and
\(\sum_\nu\operatorname{rank}(g_\nu-1)=2m+6\).
Consequently
\[
               2m+6\ge\frac{15bm}{64}
                      \ge\frac{45m}{16},
               \qquad m\le\frac{96}{13}<27.
 \tag{F1T35.16}
\]
This contradicts the original source gate and
proves that \(\beta(z_0)=1\) whenever
\(r\ge10^{31}\).

\paragraph{4. The exact remaining quotient and framework maps.}
The preceding conclusion produces the actual
surjective map
\[
 \overline\beta:\Omega_N(q)\longrightarrow L,
 \qquad \overline\beta(\phi(h))=\beta(h),
 \qquad \beta=\overline\beta\phi.
 \tag{F1T35.17}
\]
It is well defined because the displayed
two-element kernel of \(\phi\) is killed.
Equation (F1T35.7) gives an actual descended
automorphism \(\overline a_g\) of
\(\Omega_N(q)\), and the two exact squares are
\[
 \phi\widetilde a_g=\overline a_g\phi,
 \qquad \overline\beta\,\overline a_g
                  =c_g\overline\beta.
 \tag{F1T35.18}
\]
Thus the remaining orthogonal quotient is
specified together with the original outer
action, rather than chosen by a name or by
equality of character degrees. This does not
assert that this remaining quotient produces
a source.

The maps \(\pi,\beta,\phi,\overline\beta\)
and the central scalar all act on the original
coefficient representation. The original
cover, local system, inertia tuple and
middle-extension cohomology have identity
comparisons throughout. The original primitive
Jacobian corner, its integral-cleared image,
restricted polarization and cup form, Poincar\'e
biextension, Hodge filtration and period map,
and the Grothendieck--Adams realizations of
F01--F04 therefore have identity comparison
maps with pairing multiplier one. No additional
detector component is claimed.
